\documentclass{article}
\usepackage[utf8]{inputenc}
\usepackage{geometry}
\usepackage{listings}
\usepackage{xcolor}
\usepackage{hyperref}

\geometry{a4paper, margin=1in}

\title{Practical Work 7: Server-Side Sudoku on GAE}
\author{Group 16}
\date{\today}

\begin{document}

\maketitle

\section{Overview}
This practical work demonstrates the deployment of a Java-based web application on Google App Engine (GAE). Unlike static implementations, our group focused on a server-side rendering approach where the Sudoku puzzle logic is handled by a Java Servlet before being sent to the client.

\section{Architecture & Implementation}
The application follows the Model-View-Controller (MVC) pattern adapted for a simple Servlet container.

\subsection{File Structure}
\begin{itemize}
    \item \texttt{WEB-INF/appengine-web.xml}: Configures the GAE runtime environment (Java 8).
    \item \texttt{com.usth.lab7.GameController.java}: The Servlet acting as the Controller. It initializes the game matrix.
    \item \texttt{index.jsp}: The View component that dynamically renders the HTML table based on data received from the Controller.
\end{itemize}

\subsection{Key Code Snippets}
\textbf{Controller Logic (Java):}
The puzzle is defined as a 2D array and attached to the request object.
\begin{lstlisting}[language=Java, basicstyle=\ttfamily\small, breaklines=true, frame=single]
public class GameController extends HttpServlet {
    @Override
    protected void doGet(HttpServletRequest req, HttpServletResponse resp) 
    throws ServletException, IOException {
        int[][] puzzle = { {5, 3, 0...}, ... }; // Matrix initialization
        req.setAttribute("puzzle", puzzle);
        req.getRequestDispatcher("/index.jsp").forward(req, resp);
    }
}
\end{lstlisting}

\textbf{View Logic (JSP):}
We use Java scriptlets to iterate through the matrix and render inputs.
\begin{lstlisting}[language=HTML, basicstyle=\ttfamily\small, breaklines=true, frame=single]
<% for(int i=0; i<9; i++) { %>
    <tr>
    <% for(int j=0; j<9; j++) { 
         int val = puzzle[i][j]; %>
        <td>
            <input value="<%= (val!=0) ? val : "" %>" />
        </td>
    <% } %>
    </tr>
<% } %>
\end{lstlisting}

\section{Task Distribution}
The project workload was divided to ensure parallel development:

\begin{description}
    \item[Student A:] Responsible for the Java backend logic (\texttt{GameController.java}) and data structure design.
    \item[Student B:] Responsible for the frontend integration (\texttt{index.jsp}), CSS styling, and JavaScript validation hooks.
    \item[Student C:] Managed the Google Cloud Platform configuration, deployment scripts, and writing the \LaTeX\ report.
\end{description}

\section{Conclusion}
The application was successfully deployed to the Google Cloud Platform. The server-side rendering approach ensures that the initial game state is consistently managed by the backend, fulfilling the requirements of a cloud-based web application.

\end{document}